\chapter{综合题的模型入口}

\section{渐近线取两项}
\recognize{问水平、竖直或斜渐近线，函数含根式、幂指数或反三角函数。无穷远先看主阶，再取常数项；有限点先查定义域边界。}
\core{若$\lim_{x\to x_0}f(x)=\infty$（任一侧），则$x=x_0$为竖直渐近线。若$k=\lim f(x)/x$、$b=\lim[f(x)-kx]$有限，则$y=kx+b$为渐近线，$k=0$为水平线。$x\to+\infty$与$-\infty$分别检查。}
\process{\begin{enumerate}\item 查有限奇点及单侧极限。\item 无穷远比较分子分母最高阶，求斜率$k$。\item 用有理化、等价无穷小或泰勒取下一项$b$。\item 验证$f-(kx+b)\to0$。\end{enumerate}}
\fast{知道$f(x)\sim x$只能得斜率1，不能省略截距。根式先提取$x$，再用$(1+u)^\alpha=1+\alpha u+o(u)$。}
\budget{易至中；小题2--3分钟，大题6--9分钟。}
\fullquestion{math2-2025-12}
\apply{$f(x)=\sqrt[3]{x^3-3x^2+1}=x(1-3/x+1/x^3)^{1/3}$。读出主项$x$；为求截距只需展开一次，$f-x=x[-1/x+o(1/x)]\to-1$，故两端斜渐近线均为$y=x-1$。实立方根在全实轴连续，无有限竖直渐近线。}
\related{2012-01（约分后查竖直线）；2014-02（振荡余项）；2016-09；2017-09；2020-15（幂指数）；2021-18；2023-01。}

\section{极值拐点看变号}
\recognize{问极值点、拐点、凹凸区间或参数范围。极值查$f'$的符号变化，拐点查$f''$的符号变化。}
\core{$f'$由正变负为极大、负变正为极小；$f''$在两侧异号对应拐点（函数在该点连续）。$f'(x_0)=0$和$f''(x_0)=0$分别只是常见候选条件。指数因子恒正，可直接判断其余因子。}
\process{\begin{enumerate}\item 列出导数为0及不存在的候选点。\item 因式分解导数，按区间列符号表。\item 极值和拐点分别判断。\item 参数题区分单根与偶重根，保留临界参数逐一检查。\end{enumerate}}
\fast{二次多项式的双根不变号；所以“没有极值”可以允许判别式等于0，而“存在拐点”常要求判别式严格大于0。导数图题，$f'$的过零看极值，$f'$的升降转换看拐点。}
\budget{中；选择题2--3分钟，大题7--10分钟。}
\fullquestion{math2-2023-07}
\apply{$f'=(x^2+2x+a)e^x$。无极值要求二次式没有变号的实根，即$4-4a\le0$，得$a\ge1$。$f''=(x^2+4x+a+2)e^x$；存在拐点要求有两个单根，$16-4(a+2)>0$，得$a<2$。取交集$[1,2)$，选C。}
\related{2015-04、2016-04（导数图）；2016-05（曲率与相切）；2018-10；2019-02；2025-02（变上限积分局部结构）。}

\section{切线法线先定点}
\recognize{隐式、参数式、极坐标曲线给一个坐标或参数，求切线、法线或截距。先找准确曲线点，再求斜率。}
\core{切线$Y-y_0=m(X-x_0)$；$m\ne0$时法线斜率$-1/m$。参数曲线$m=y_t'/x_t'$；极坐标先写$x=r\cos\theta,y=r\sin\theta$。若$x_t'=0,y_t'\ne0$，切线竖直。}
\process{\begin{enumerate}\item 用给定坐标求另一坐标或对应参数。\item 对原关系求导，代点得到切线方向。\item 求法线时再取负倒数。\item 截距由直线方程令$X=0$或$Y=0$。\end{enumerate}}
\fast{切线水平时法线竖直，不能硬写无限大的斜率。参数导数同时为0时，用邻近展开或$dy/dx$的极限判断方向。}
\budget{易至中；2--3分钟。}
\fullquestion{math2-2023-14}
\apply{$x=1$代原方程得$y^5+2y^3=3$，其左侧严格递增，故唯一对应$y=1$。隐式求导$9x^2=(5y^4+6y^2)y'$，代点得切线斜率$9/11$；题目要法线，答案为$-11/9$。}
\related{2013-12；2014-12（极坐标切线）；2019-10（参数切线截距）；2021-20（截距再取最值）。}

\section{曲率圆找法向}
\recognize{问曲率、曲率半径、曲率圆或同切线曲线的弯曲程度。先选在指定点不退化的曲线表示。}
\core{$\kappa=|y''|/[1+(y')^2]^{3/2}$，$R=1/\kappa$。参数式$\kappa=|x'y''-y'x''|/[(x')^2+(y')^2]^{3/2}$。若$y''\ne0$，曲率圆心$\left(x-\frac{y'(1+(y')^2)}{y''},\ y+\frac{1+(y')^2}{y''}\right)$。}
\process{\begin{enumerate}\item 选择$y(x)$、$x(y)$或参数表示。\item 算一、二阶导并求曲率。\item 取倒数求半径。\item 圆心沿法线指向局部弯曲一侧，写圆方程。\end{enumerate}}
\fast{竖直切线不适合强用$y(x)$公式；交换$x,y$或使用参数式即可。曲率为0时没有有限曲率圆。}
\budget{中；3--4分钟。}
\fullquestion{math2-2024-11}
\apply{原点处$y(x)$有竖直切线，改看$x=y^2$。在$y=0$时$x_y'=0,x_{yy}''=2$，故$\kappa=2,R=1/2$。曲线向右开口，圆心沿正$x$方向，故为$(1/2,0)$，曲率圆$(x-1/2)^2+y^2=1/4$。}
\related{2012-13；2014-04、2018-12（参数曲率）；2019-06（相切与等曲率）；2026-13（隐式二阶导求半径）。}

\section{零点与参数}
\recognize{问零点个数、恰有两根或参数范围。先把方程写成$F(x)=0$，圈出定义域、端点极限和参数所在位置。}
\core{连续函数异号保证有根；严格单调保证至多一根。若$F'$仅一处变号，先找极值再比较极值与0。将参数分离为$a=\phi(x)$时，根数等于水平线与$\phi$图形的交点数。}
\process{\begin{enumerate}\item 写定义域和两端极限。\item 求导并列单调表；参数优先从导数中分离。\item 用极值符号判根数，等号对应重根。\item 每个单调区间分别验证存在性。\end{enumerate}}
\fast{“存在”与“唯一”分两步落笔。原式含难积函数时先求导，避免为了找根强行求原函数。}
\budget{中；选择题2--3分钟，证明或参数讨论8--12分钟。}
\fullquestion{math2-2015-19}
\apply{变上限积分无需先算闭式。直接得$f'(x)=\sqrt{1+x^2}(2x-1)$，所以在$(-\infty,1/2)$递减、$(1/2,\infty)$递增。积分限交会给$f(1)=0$，严格递增说明$f(1/2)<0$；又$f(x)\to+\infty$（$x\to-\infty$），故左支一根、右支一根，共两根。}
\related{2012-21（唯一根接数列）；2016-21（原函数零点）；2021-04（参数与两根）。}

\section{递推先证收敛}
\recognize{给$x_{n+1}=\phi(x_n)$或隐式递推，问极限、收敛或无穷小比较。先证明迭代始终留在哪个区间。}
\core{单调有界数列必收敛。只有先确认收敛、且$\phi$在极限处连续，才能令$L=\phi(L)$。常用$e^x>1+x$、$\sin x<x$（$x>0$的适用小区间）。}
\process{\begin{enumerate}\item 归纳证明正性与不变区间。\item 比较$\phi(x)$与$x$，得到单调性。\item 由单调有界得到极限存在。\item 代入递推并证明候选极限唯一。\end{enumerate}}
\fast{隐式指数递推先取对数；比较$x_{n+1}-x_n$时可合并成一个对数。两递推比较速度，研究比值递推比逐项求通项更短。}
\budget{中至难；8--12分钟，选择题2--3分钟。}
\fullquestion{math2-2018-21}
\apply{原式给$x_{n+1}=\ln[(e^{x_n}-1)/x_n]$。对$x>0$，$e^x-1>x$保证下一项为正；$e^x-1<xe^x$保证下一项小于$x$。故递减有下界。极限满足$Le^L=e^L-1$，而$1-e^{-L}<L$对$L>0$严格成立，故只能$L=0$。}
\related{2012-21；2023-03（正弦迭代与平方迭代比较）；2024-04（复合数列的逆向推断）。}

\section{一阶项凑导数}
\recognize{条件或结论出现$f'+af$、$xf'-f$、函数值之比与指数。先找乘积或商的导数。}
\core{$(e^{ax}f)'=e^{ax}(f'+af)$；$(f/x)'=(xf'-f)/x^2$；$f>0$时$(\ln f)'=f'/f$。}
\process{\begin{enumerate}\item 把目标移到一侧，识别导数线性组合。\item 乘积分因子或除以正函数，构造$G$。\item 用$G'$的符号或罗尔定理。\item 最后还原为原函数的大小或零点。\end{enumerate}}
\fast{看见$f'>f>0$，立即想到$e^{-x}f$递增或$\ln f-x$递增；直接多次求$f$的导数没有额外信息。}
\budget{中；选择题1.5--2.5分钟，证明6--10分钟。}
\fullquestion{math2-2020-06}
\apply{$G=e^{-x}f$满足$G'=e^{-x}(f'-f)>0$。任意$u<v$都有$f(v)/f(u)>e^{v-u}$。把各选项两点顺序统一后比较即可，指数恰好来自积分因子。}
\related{2013-18（$f''+f'$）；2017-19（组合导数零点）；2015-21（切线与凸性）。}

\section{端点插值与余项}
\recognize{出现$f(a),f(b)$、线性插值、$x(1-x)$及二阶导数界。先减去端点连线，使新函数两端为零。}
\core{设$L(x)=\frac{b-x}{b-a}f(a)+\frac{x-a}{b-a}f(b)$。若$|f''|\le M$，则$|f(x)-L(x)|\le\frac M2(x-a)(b-x)$；积分后梯形误差绝对值不超过$M(b-a)^3/12$。}
\process{\begin{enumerate}\item 构造$F=f-L$，得$F(a)=F(b)=0$。\item 再减去待证上界的二次函数。\item 用二阶导数符号判断图形在弦线哪侧。\item 下界同理，积分得到整体误差。\end{enumerate}}
\fast{目标中的$x(1-x)$两端为0、二阶导数为常数，正好承接$|f''|$；分别对两端做泰勒会引入两个不同中间点，常使证明更长。}
\budget{难；10--14分钟。}
\fullquestion{math2-2024-21}
\apply{令$F=f-f(0)(1-x)-f(1)x$。$F-\frac12x(1-x)$二阶导数$f''+1\ge0$，端点为0，故不大于0；$F+\frac12x(1-x)$二阶导数$f''-1\le0$，故不小于0。积分上界$\frac12\int_0^1x(1-x)dx=1/12$直接给第二问。}
\related{2014-03（弦线比较）；2023-21（两端泰勒与二阶导数）；2022-21（中点积分不等式）。}

\section{割线斜率与凸性}
\recognize{问$f'$单调、三点割线斜率比较，或中点函数值与积分均值比较。把“导数增长”翻译成“后段斜率更大”。}
\core{若$f'$严格递增，则$x_1<x_2<x_3$时，$\frac{f(x_2)-f(x_1)}{x_2-x_1}<\frac{f(x_3)-f(x_2)}{x_3-x_2}$。$f''\ge0$时，$f((a+b)/2)\le\frac1{b-a}\int_a^bf$。}
\process{\begin{enumerate}\item 正向用相邻两区间的拉格朗日定理。\item 逆向固定两点，把第三点靠近端点，获得导数与割线关系。\item 先得$f'$非减，再用严格割线条件排除常值区间。\item 积分均值型用关于中点的对称配对。\end{enumerate}}
\fast{“严格递增”不等于处处$f''>0$；例如$x^4$的导数严格递增但二阶导数在0为0。逆向证明的严格性要靠排除线性区间。}
\budget{难；10--15分钟。}
\fullquestion{math2-2025-21}
\apply{正向在$(x_1,x_2)$、$(x_2,x_3)$各取中值点，前者导数严格小于后者。逆向令割线端点趋近，先得$f'(u)\le f'(v)$；若某$u<v$导数相等，非减性使$f'$在$[u,v]$恒定，$f$在该段线性，与任意内部三点的严格斜率不等式冲突。}
\related{2022-21；2026-05（固定端点斜率与整体凸性的区别）。}

\section{截面与旋转体}
\recognize{问面积、绕轴体积、由体积比反求参数。先画截面，标明外半径、内半径和厚度。}
\core{$S=\int_a^b|f-g|dx$；垫片法$V=\pi\int(R^2-r^2)dx$；圆柱壳法$V=2\pi\int \rho h\,d\rho$。旋转轴平移时半径是到轴的距离。}
\process{\begin{enumerate}\item 求交点、画区域。\item 比较横切和竖切，选择少分段的方向。\item 写半径与厚度，再列积分。\item 用长度三次方核查体积量纲。\end{enumerate}}
\fast{绕$y$轴且区域写成$0\le y\le f(x)$，壳法通常避免求反函数；旋转轴穿过区域时，先确定实际截面，不能重复计算重叠部分。}
\budget{中；8--12分钟。}
\fullquestion{math2-2015-16}
\apply{区域$0\le x\le\pi/2$、$0\le y\le A\sin x$。绕$x$轴用垫片得$V_1=\pi A^2\int_0^{\pi/2}\sin^2x\,dx=\pi^2A^2/4$；绕$y$轴用壳法得$V_2=2\pi A\int_0^{\pi/2}x\sin x\,dx=2\pi A$。由$A>0$约去一次$A$，得$A=8/\pi$。}
\related{2012-17；2013-16；2014-21（平移轴）；2020-18（以$y$作积分变量）。}

\section{无界区域与变限体积}
\recognize{区域延伸到无穷、边界含参数$t$，或先求拐点再求体积。几何积分与广义积分、极值会串在一起。}
\core{无界体积$V=\pi\int_a^{\infty}R^2dx$需验证收敛。$V(t)=\pi\int_{\alpha(t)}^{\beta(t)}h(x)dx$满足$V'=\pi[h(\beta)\beta'-h(\alpha)\alpha']$。}
\process{\begin{enumerate}\item 先求确定边界的几何点。\item 有界段与尾部拆开列式。\item 参数极值题先对积分求导，避免先求复杂原函数。\item 比较驻点与参数边界的极限。\end{enumerate}}
\fast{无界区域面积与旋转体积分别比较$f$与$f^2$，一个发散不决定另一个。}
\budget{中至难；10--14分钟。}
\fullquestion{math2-2024-19}
\apply{$y=\sqrt{x}e^{-x}$，故$V(t)=\pi\int_t^{2t}xe^{-2x}dx$。直接求导得$V'=\pi t e^{-4t}(4-e^{2t})$，以$t=\ln2$为分界先增后减，得到最大体积；这比先展开$V(t)$再求导更短。}
\related{2023-19（无界面积与体积）；2026-20（拐点、直线段、无界尾部）。}

\section{弧长与曲面面积}
\recognize{给曲线、参数方程或由积分确定的函数，求长度、旋转曲面面积。先算$\sqrt{1+(y')^2}$是否能化为完全平方。}
\core{$s=\int\sqrt{1+(y')^2}dx$；参数式$s=\int\sqrt{(x')^2+(y')^2}dt$；绕$x$轴曲面面积$A=2\pi\int |y|\,ds$。}
\process{\begin{enumerate}\item 由题给关系求出$y'$。\item 配方化简弧长因子，处理绝对值。\item 长度积分一次；面积再乘$2\pi$与旋转半径。\item 参数区间若有三角函数变号，按符号分段。\end{enumerate}}
\fast{若$y'=u-1/(4u)$，则$1+(y')^2=(u+1/(4u))^2$。看到二项互为倒数，先试完全平方。}
\budget{中；填空2--3分钟，大题8--12分钟。}
\fullquestion{math2-2022-18}
\apply{先将方程化为$y'-2y/x=(2\ln x-1)/(2x)$，乘$x^{-2}$，有$(y/x^2)'=(\ln x-1/2)/x^3$。积分并用$y(1)=1/4$得$y=x^2/4-(\ln x)/2$。于是$y'=x/2-1/(2x)$，弧长因子恰为$x/2+1/(2x)$，从1积至$e$，弧长$(e^2+1)/4$。}
\related{2013-21；2016-20（参数曲面）；2019-12；2021-19；2023-12。}

\section{物理量取投影}
\recognize{引力、功、水压力、质心或变化率。先确定所求量的方向、微元和单位，再写积分。}
\core{引力微元$dF=Gm\,dm/r^2$；沿位移方向做功$dW=F_{\parallel}dx$；静水压力$dP=\rho gh\,dA$；质心$\bar x=\frac{\int x\,dm}{\int dm}$。均匀薄片$g(x)\le y\le f(x)$的$\bar x=\int x(f-g)dx/S$、$\bar y=\int(f^2-g^2)dx/(2S)$；非均匀细棒$dm=\rho(x)dx$。平均值$\bar f=\frac1{b-a}\int_a^bf$。}
\process{\begin{enumerate}\item 选位置变量，画距离和夹角。\item 写质量/面积微元。\item 取受力或运动方向分量，对称分量先抵消。\item 积分并核查单位。\end{enumerate}}
\fast{引力分母由$r^2$再乘方向余弦中的$r$，常变成$r^3$；遗漏投影会少一层根号。变化率题先写$V(r,h)$再对时间整体求导。}
\budget{中；选择题2--3分钟，建模大题8--12分钟。}
\fullquestion{math2-2025-06}
\apply{位置$(x,0)$到$Q$距离$r=\sqrt{1+x^2}$，引力沿负$x$方向的分量大小为$Gx/r^3$。克服引力的功$W=\int_0^lGx/(1+x^2)^{3/2}dx$，对应B。分子$x$来自方向投影。}
\fullquestion{math2-2026-04}
\apply{棒上微元$dm=dx$，左右水平分量抵消；竖直投影因子为$1/r$。因此总力$F=2Gm\int_0^1(1+x^2)^{-3/2}dx$，对应B。与上一题同为$r^{-3}$，分子不同源于投影方向。}
\related{2014-13（质心）；2016-13、2018-20、2021-03（相关变化率）；2020-12（水压力）；2024-15、2026-15（平均值）。}

\section{几何条件建方程}
\recognize{曲线未知，给切线截距、法线截距、面积恒比。将几何量写成$x,y,y'$，再化成微分方程。}
\core{切线$Y-y=y'(X-x)$；纵截距$b=y-xy'$；横截距$a=x-y/y'$（$y'\ne0$）；法线斜率$-1/y'$。面积函数$S(x)=\int_0^xy(t)dt$满足$S'=y$。}
\process{\begin{enumerate}\item 设一般点$(x,y)$，写切线/法线。\item 求截距或三角形面积。\item 把恒等几何条件翻译成方程。\item 积分条件先求导降为微分方程，再代初值。\end{enumerate}}
\fast{截距可以有符号，距离和面积必须取绝对值；先读题限定象限。含$y-xy'$的条件适合研究$(y/x)'$。}
\budget{中至难；10--14分钟。}
\fullquestion{math2-2020-21}
\apply{切线与$x$轴交点横坐标$x-y/y'$，故三角形面积$y^2/(2y')$。由面积比得$\int_0^xy(t)dt=3y^2/(4y')$。求导整理得$2(y')^2=3yy''$，因$x>0$时$y,y'>0$，可积分为$\ln y'=\frac23\ln y+C$，结合过原点得$y=kx^3$（$k>0$）。读题中“面积比恒定”提示先引入面积函数。}
\related{2017-21；2023-17（截距建一阶线性方程再求最值）；2021-20（法线截距最小）。}

\section{一阶线性与变量互换}
\recognize{方程能写成$y'+P(x)y=Q(x)$；若关于$y$非线性，检查能否改为$x=x(y)$后线性。}
\core{积分因子$\mu=e^{\int Pdx}$，$(\mu y)'=\mu Q$，$y=\mu^{-1}(\int\mu Qdx+C)$。自变量互换时$dx/dy=1/y'$应在合法区间使用。}
\process{\begin{enumerate}\item 标准化并读出$P,Q$。\item 求积分因子，写成完整乘积导数。\item 积分后代初值。\item 核查定义区间和原方程。\end{enumerate}}
\fast{形如$y\,dx+(x-3y^2)dy=0$，把$x$当未知函数立即线性；硬把$dy/dx$分离会把$x,y$缠在分母。}
\budget{易至中；填空2--4分钟，大题第一问4--6分钟。}
\fullquestion{math2-2012-12}
\apply{改写$dx/dy+x/y=3y$。积分因子为$y$，故$(xy)'=3y^2$，得到$xy=y^3+C$。初值$(x,y)=(1,1)$给$C=0$，于是初值附近$y=\sqrt{x}$，正根由初值确定。}
\related{2019-17（接体积）；2021-20（接截距）；2022-18（接弧长）。}

\section{比例定律建模}
\recognize{变化率与现存量、温差或混合浓度成比例。把“变化率”翻译为导数，先确定正负方向。}
\core{牛顿冷却定律$T'=-k(T-T_a)$（$k>0$），解为$T-T_a=(T_0-T_a)e^{-kt}$。混合问题满足“溶质量变化率＝流入速率乘流入浓度－流出速率乘当前浓度”。}
\process{\begin{enumerate}\item 设未知量及单位，写变化率平衡。\item 用初始条件确定初值。\item 用第二个观测值求比例系数。\item 计算剩余时间时统一起算时刻或取温差比。\end{enumerate}}
\fast{冷却衰减的是相对环境的温差，不是温度本身；同一环境下每缩小10倍温差需要相同时间。}
\budget{易至中；6--9分钟。}
\fullquestion{math2-2015-20}
\apply{设环境温度20，初始温差100，30分钟后温差10，故$e^{-30k}=1/10$。从30降到21是温差从10降到1，也缩小10倍，所以还需30分钟。先认出相同温差比可省去显式求$k$。}
\related{2021-03（体积与表面积相关变化率）；2017-06（速度图面积与相对位移）。}

\section{换元化为可分离}
\recognize{右侧只含$x+y$、$y/x$或线性组合。把重复结构整体设为新变量。}
\core{$y'=F(x+y)$令$u=x+y$，得$u'=1+F(u)$；齐次方程$y'=F(y/x)$令$y=ux$，得$x u'=F(u)-u$。}
\process{\begin{enumerate}\item 找反复出现的组合。\item 换元同时写出导数关系。\item 分离变量积分。\item 代初值，并检查除去的因子是否对应常值解。\end{enumerate}}
\fast{先看结构再决定方程名称。含$dx,dy$的二次齐次式也可检查是否恰为某二次式的全微分。}
\budget{中；填空3--4分钟，大题6--8分钟。}
\fullquestion{math2-2024-13}
\apply{令$u=x+y$，得$u'=1+u^{-2}$，故$dx=u^2/(1+u^2)\,du$。积分得到$x=u-\arctan u+C$，代$u(1)=1$后可写$y=\arctan(x+y)-\pi/4$。初值处$u=1$使除法合法。}
\related{2014-16（隐式解接极值）；2025-15（恰当微分与齐次换元）。}

\section{特征根与解的行为}
\recognize{常系数线性微分方程求通解，或要求任意解有界、广义积分均收敛。把条件翻译成每一个特征模态的行为。}
\core{特征根$r$为$k$重，贡献$e^{rx}(C_0+\cdots+C_{k-1}x^{k-1})$；复根$\alpha\pm i\beta$对应$e^{\alpha x}(C\cos\beta x+D\sin\beta x)$。所有模态在$[0,\infty)$可积，需实部严格小于0。}
\process{\begin{enumerate}\item 写特征多项式。\item 分根型列通解。\item “任意解”逐模态检查，不能靠系数抵消。\item 二次式$r^2+pr+q$全根实部为负等价$p>0,q>0$。\end{enumerate}}
\fast{整个实轴有界比正半轴衰减更强：实指数在某一端必增长；纯虚根可有界，但重纯虚根的$x$因子会破坏有界。}
\budget{中；选择题2--3分钟，求通解填空2--4分钟。}
\fullquestion{math2-2025-03}
\apply{特征多项式$r^2-2ar+a+2$，对照$p=-2a,q=a+2$，得$a<0$且$a>-2$，即$(-2,0)$。端点$a=0$的纯振荡积分一般没有极限，$a=-2$有常数模态，都须排除。}
\related{2021-15、2022-14（三阶通解）；2023-04（全实轴有界）；2020-13（解后积分）。}

\section{特解叠加与待定系数}
\recognize{给多个非齐次特解，问线性组合；或右端为指数乘多项式/三角函数。先分清齐次与非齐次。}
\core{线性算子$L$满足$L[y_i]=r$时，$L[\sum c_i y_i]=(\sum c_i)r$。同一非齐次方程两特解之差为齐次解。待定特解$e^{\alpha x}P_m(x)$遇$\alpha$为$k$重特征根时乘$x^k$。}
\process{\begin{enumerate}\item 标记每个已知解的右端。\item 线性组合只比较系数和。\item 求方程时先由解之差提取特征根。\item 最后用任一特解反推右端或待定常数。\end{enumerate}}
\fast{通解中的任意常数部分对应齐次方程；固定项才是特解。特征根与右端指数重合时必须乘$x$，重合两重则乘$x^2$。}
\budget{易至中；选择题1--2分钟，填空2--4分钟。}
\fullquestion{math2-2026-02}
\apply{把方程写作$L[y]=r\ne0$。第一组合仍为非齐次解，得$2\lambda+\mu=1$；第二组合为齐次解，得$\lambda-2\mu=0$，解得$\lambda=2/5,\mu=1/5$，对应B。全程无需知道$y_1,y_2$的具体表达式。}
\related{2013-13；2016-11；2019-04；2017-04（网站题干缺失，见附件）。}

\section{变系数降阶}
\recognize{欧拉型$x^2y''+axy'+by$、方程不显含$y$、或已知一个齐次解。分别试$x=e^t$、$p=y'$、$y=u y_1$。}
\core{欧拉变换给$xy'=\dot y$、$x^2y''=\ddot y-\dot y$。缺$y$时$p=y'$降一阶；若$p'=a(x)p+b(x)p^2$，再令$u=1/p$化线性。}
\process{\begin{enumerate}\item 先判断缺变量和系数幂次。\item 换元只降阶，不同时展开所有积分。\item 用导数初值先求$p$，再积分求$y$。\item 除以$p$前检查$p=0$是否另为解。\end{enumerate}}
\fast{缺$y$的非线性二阶方程通常没有必要寻找二阶通解；已知$y'(x_0)\ne0$可在初值邻域使用倒数换元。}
\budget{中至难；10--14分钟。}
\fullquestion{math2-2026-21}
\apply{缺$y$提示令$p=y'$，得$x^2p'-2xp-p^2=0$；再令$u=1/p$，化为$(x^2u)'=-1$。由$p(3)=-9$得$p=x^2/(2-x)$；长除法后$p=-x-2-4/(x-2)$。再用$y(3)=1/2$得$y=-x^2/2-2x-4\ln(x-2)+11$。}
\related{2016-19（已知解令$y=u e^x$）；2024-18（欧拉变换后求积分）。}

\section{偏导与可微分层}
\recognize{分段二元函数、振荡因子、原点处偏导和可微。分别沿坐标轴求偏导，再检验二维余项。}
\core{可微要求$f(h,k)-f(0,0)-Ah-Bk=o(\sqrt{h^2+k^2})$；偏导存在只确定$A,B$，不能替代此极限。连续偏导是可微的充分条件。}
\process{\begin{enumerate}\item 由定义求原点两偏导。\item 写线性主部$Ah+Bk$。\item 用$r=\sqrt{h^2+k^2}$统一估计余项。\item 偏导连续性另求邻点偏导，并找趋近路径。\end{enumerate}}
\fast{振荡有界时，$O(r^2)$余项除$r$即趋0；原点可微仍可能邻点偏导无界。路径法能证不存在，不能仅靠若干路径证明存在。}
\budget{中；选择题2--3分钟，证明5--8分钟。}
\fullquestion{math2-2024-05}
\apply{坐标轴上函数为0，两偏导为0。且$|f(h,k)|\le h^2+k^2=r^2$，故$|f|/r\le r\to0$，原点可微。$xy\ne0$时$f_x=2x\sin(1/(xy))-(x^2+y^2)\cos(1/(xy))/(x^2y)$，取适当$x=y\to0$使余弦为1，第二项不趋0，偏导不连续，选C。}
\related{2020-05（累次极限与二元极限）；2012-05、2017-05（分变量单调性）。}

\section{链式法则逐层展开}
\recognize{$z=f(u,v)$，$u,v$又由$x,y$确定，问偏导组合、二阶导或由两条路径求全微分。先画依赖层次再写算子。}
\core{$\partial_x=u_x\partial_u+v_x\partial_v$；$z_x=f_u u_x+f_v v_x$。线性内变量时，$z_{xx}=f_{uu}u_x^2+2f_{uv}u_xv_x+f_{vv}v_x^2$；非线性另加$f_u u_{xx}+f_v v_{xx}$。}
\process{\begin{enumerate}\item 明确独立变量与内变量。\item 一阶写链式法则。\item 二阶同时求外层和内层导数。\item 最后再代指定点。\end{enumerate}}
\fast{给偏导组合时，先组合微分算子再展开，往往大幅抵消。两条路径穿过同一点时，对两恒等式各求导，得到$f_x,f_y$的二元方程组。}
\budget{中；小题2--3分钟，大题8--12分钟。}
\fullquestion{math2-2024-20}
\apply{由$u=2x+y,v=3x-y$知$\partial_x=2\partial_u+3\partial_v$、$\partial_y=\partial_u-\partial_v$。目标算子$(\partial_x+3\partial_y)(\partial_x-2\partial_y)=25\partial_u\partial_v$，立即得$f_{uv}=1/25$。再按给定轴上数据两次积分，得$f=uv/25-(u+1)e^{-u}+v^2/50$。}
\related{2017-16；2019-20（指数因子消一阶导）；2021-06（两路径求全微分）；2022-04；2026-14。}

\section{隐函数偏导}
\recognize{$F(x,y,z)=0$确定$z(x,y)$，问偏导、全微分或二阶偏导。先求指定点的$z$值。}
\core{$F_z\ne0$时$z_x=-F_x/F_z$、$z_y=-F_y/F_z$，$dz=-(F_xdx+F_ydy)/F_z$；$F_{xx}+2F_{xz}z_x+F_{zz}z_x^2+F_z z_{xx}=0$。}
\process{\begin{enumerate}\item 代入$(x_0,y_0)$解$z_0$。\item 写一阶隐式求导关系。\item 二阶再对整个关系求导。\item 代点后解未知导数。\end{enumerate}}
\fast{求单点二阶偏导时，在一阶恒等式求导后代点即可，无需先整理出完整$z_x$的分式。变上限积分只对相应上下限求导。}
\budget{易至中；2--4分钟。}
\fullquestion{math2-2025-01}
\apply{移项为$z+\ln z=\int_y^xe^{-t^2}dt$。分别求偏导得$(1+1/z)z_x=e^{-x^2}$、$(1+1/z)z_y=-e^{-y^2}$。相加即$z_x+z_y=\frac z{z+1}(e^{-x^2}-e^{-y^2})$，选A。下限的负号决定选项。}
\related{2014-11；2015-13；2018-13；2021-13；2023-13；2026-03。}

\section{从微分还原函数}
\recognize{给$df=Pdx+Qdy$、混合偏导和边界值，求$f$再求极值。“积分常数”要保留为另一变量的函数。}
\core{由$f_x=P$积分得$f=\int Pdx+\phi(y)$，再与$f_y=Q$比较求$\phi'$。给$f_{xy}$时逐次积分，通常产生两个一元函数。}
\process{\begin{enumerate}\item 选最易积分的一阶偏导。\item 对该变量积分并保留未知函数。\item 用另一偏导及边界值确定未知函数。\item 代回两偏导检查后再求驻点。\end{enumerate}}
\fast{边界$f(x,0)$适合确定只含$x$的函数；不要在偏积分时把另一变量函数写成常数$C$。}
\budget{中；8--12分钟。}
\fullquestion{math2-2025-19}
\apply{由$f_x=-2xe^{-y}$得$f=-x^2e^{-y}+\phi(y)$。对$y$求导比较可得$\phi'=-(y+1)e^{-y}$，故$f=(y+2-x^2)e^{-y}$，初值消掉常数。驻点$(0,-1)$处$A=-2e,B=0,C=-e$，得极大值$e$。}
\related{2015-17；2017-12；2022-20（先换元还原函数）。}

\section{二元极值与退化点}
\recognize{求$f(x,y)$极值、极值点或由隐函数确定的极值。先解两个一阶偏导同时为0。}
\core{驻点处$A=f_{xx},B=f_{xy},C=f_{yy}$。$AC-B^2>0$时$A>0$极小、$A<0$极大；$AC-B^2<0$为鞍点；等于0时回到增量符号、限制路径或高阶项。}
\process{\begin{enumerate}\item 提取永不为0的指数因子。\item 用两方程相减、相加或代入求全体驻点。\item 计算三项二阶偏导并判别。\item 退化点比较邻域增量，区分局部与全局。\end{enumerate}}
\fast{先由较简单的$f_y=0$解$y$，再代$f_x=0$；不必一开始就求完整Hessian。对隐函数极值，驻点处$z_x=z_y=0$可大量删去二阶求导项。}
\budget{中；填空3--4分钟，大题8--12分钟。}
\fullquestion{math2-2026-19}
\apply{$f_y=-2ye^x$先给$y=0$，再由$f_x=e^x(2x^2+4x-y^2)=0$得$(0,0),(-2,0)$。原点判别式$-8<0$是鞍点；$(-2,0)$处$A=-4e^{-2},B=0,C=-2e^{-2}$，为极大值点，极大值$8e^{-2}$。}
\related{2012-16；2016-17（隐函数极值）；2020-17；2023-18（周期驻点）；2024-12（可分离结构）。}

\section{约束最值与建模}
\recognize{限定曲线、资源总量、分段长度，问最短最长或面积最小。先把目标写成最少变量的函数。}
\core{约束$g=0$且$\nabla g\ne0$时，候选点满足$\nabla f=\lambda\nabla g$。有界闭域连续函数最值还需比较边界与角点。正数和的平方模型可用柯西不等式$\sum a_i x_i^2\ge(\sum x_i)^2/(\sum1/a_i)$。}
\process{\begin{enumerate}\item 标记每个变量的实际范围。\item 能消元则消元；否则拉格朗日乘子。\item 求内部候选点。\item 检查边界是否允许并比较目标值。\end{enumerate}}
\fast{几段周长之和固定、面积均为“系数乘长度平方”时，用加权平方模型直接找最小值，比繁复求偏导更稳。}
\budget{中；8--12分钟。}
\fullquestion{math2-2018-19}
\apply{设三段长度$l_1,l_2,l_3>0$、和为2，面积和为$l_1^2/(4\pi)+l_2^2/16+l_3^2/(12\sqrt3)$。柯西给下界$4/(4\pi+16+12\sqrt3)$，等号要求三段长度比$4\pi:16:12\sqrt3$，均为正，因而下界可以达到。}
\related{2013-19（曲线约束距离）；2013-21（形心）；2022-22（线性代数中的商最值）。}

\section{换序先画投影}
\recognize{内积分无法初等求出，或问二次积分换序。将不等式还原成区域，再从另一方向切片。}
\core{$D=\{a\le x\le b,\phi_1(x)\le y\le\phi_2(x)\}$。换序时先定$y$投影，再解每条水平线与区域的交段；出现两段需分别积分。}
\process{\begin{enumerate}\item 写出全部边界并求交点。\item 画原区域，标投影范围。\item 用水平线/竖直线扫描，寻找分段点。\item 先写新积分限，再抄被积函数。\end{enumerate}}
\fast{被积函数任意时，对称区域不自动允许把左半写成右半的两倍；只有函数也相应对称才可。}
\budget{中；选择题2--3分钟，计算6--9分钟。}
\fullquestion{math2-2025-05}
\apply{区域为$-2\le x\le2$、$4-x^2\le y\le4$。投影$0\le y\le4$，水平截线分成$[-2,-\sqrt{4-y}]$与$[\sqrt{4-y},2]$两段，故选A。任意$f$不保证关于$y$轴对称，不能选择“右段乘2”。}
\related{2017-13；2020-10；2022-02；2024-06（反三角边界）。}

\section{对称先消项}
\recognize{区域关于坐标轴或$y=x$对称，被积函数含奇次项、$x-y$、轮换形式。先判断能消掉什么，再计算。}
\core{$D$关于$y$轴对称时，关于$x$奇的函数积分为0；$D$关于$y=x$对称时，$\iint_Df(x,y)=\iint_Df(y,x)$，可取两式平均。}
\process{\begin{enumerate}\item 验证整个区域在变换下不变。\item 检查被积函数变号还是不变。\item 消奇项、轮换平均，或对半区域乘2。\item 对剩余积分选坐标。\end{enumerate}}
\fast{先化简再换元。$x-y$可因轮换反对称消掉，$(x-y)^2$则保留；对称的是函数值，不能消平方中的$xy$。}
\budget{易至中；小题1--2分钟，大题6--10分钟。}
\fullquestion{math2-2024-17}
\apply{边界$xy\in[1/3,3]$、$y/x\in[1/3,3]$在交换$x,y$后不变，所以$\iint_D(x-y)=0$，原积分仅为面积。取$u=xy,v=y/x$，Jacobian绝对值$1/(2v)$，面积$\int_{1/3}^3du\int_{1/3}^3dv/(2v)=\frac83\ln3$。}
\related{2013-06；2016-18；2018-06；2026-07（三角域与正方形）。}

\section{极坐标按射线取小界}
\recognize{含$x^2+y^2$、圆、环、扇形，或区域由两圆相交。先从原点发射射线，确定先碰到哪条边界。}
\core{$x=r\cos\theta,y=r\sin\theta,dxdy=rdrd\theta$。对同一射线上的两个区间$0\le r\le R_1$、$0\le r\le R_2$，交集上界为$\min(R_1,R_2)$，并集上界为$\max(R_1,R_2)$。}
\process{\begin{enumerate}\item 画出圆心与交点，确定角域。\item 将全部边界转为$r$的条件。\item 解$R_1=R_2$找角度分界。\item 连同Jacobian一起改写被积函数。\end{enumerate}}
\fast{交集的径向上界必须满足所有不等式。先在$\theta=0$、$\pi/4$检验，能快速排除取错圆的情况。}
\budget{中至难；10--14分钟。}
\fullquestion{math2-2025-20}
\apply{两圆分别给$r\le4\cos\theta$与$r\le4\sin\theta$。轮换对称后取$0\le\theta\le\pi/4$的两倍，这一半上界必须是$4\sin\theta$。因此$I=2\int_0^{\pi/4}\int_0^{4\sin\theta}r^3(\cos\theta-\sin\theta)^2drd\theta=12\pi-112/3$。用$I\ge0$及窄角附近半径趋0检查设限。}
\related{2012-18；2014-17；2017-20；2019-18；2021-21；2023-20。}

\section{直角与极坐标比较}
\recognize{被积函数有$y\,F(\sqrt{x^2+y^2})$，同时边界为直线与圆。可比较先对$y$换元与直接极坐标两条路。}
\core{固定$x$令$u=\sqrt{x^2+y^2}$，有$y\,dy=u\,du$；极坐标则$y\,dxdy=r^2\sin\theta\,drd\theta$。}
\process{\begin{enumerate}\item 先查对称能否减半。\item 判断新变量下边界是否简化。\item 若变为三角区域，换序可能进一步分离。\item 若原域本就是扇形，直接极坐标通常最短。\end{enumerate}}
\fast{不凭根式就机械用极坐标；看内层是否已有完整微分$y\,dy$。两种路径都应先核对区域。}
\budget{中；8--12分钟。}
\fullquestion{math2-2026-17}
\apply{区域由$y\ge|x|$与$x^2+y^2\le2$组成，是$\pi/4\le\theta\le3\pi/4$的扇形。直接极坐标给$I=(\int_{\pi/4}^{3\pi/4}\sin\theta d\theta)\int_0^{\sqrt2}r^2\sin r\,dr=4\sin\sqrt2-2\sqrt2$。另一条路线令$u=\sqrt{x^2+y^2}$后换序也可，直接极坐标少一次区域重画。}
\related{2020-19（圆根式与$x$比例）；2022-19（直线与圆的角域分段）。}
